\documentclass[11pt,a4paper]{article}

% ============================================================
%  QIMMA -- COURS : Calcul intégral
%  2ème Bac Sciences Mathématiques (A et B)
%  Standard exhaustif : définitions formelles + démonstrations
%  Basé sur templates/qimma-template.tex (charte Qimma)
% ============================================================

\usepackage[french]{babel}
\usepackage[utf8]{inputenc}
\usepackage[T1]{fontenc}
\usepackage{lmodern}
\usepackage[a4paper,margin=2cm,top=2.3cm,bottom=2.6cm]{geometry}
\usepackage{amsmath,amssymb}
\usepackage{xcolor}
\usepackage{graphicx}
\usepackage{tikz}
\usetikzlibrary{shadows,calc}
\usepackage[most]{tcolorbox}
\usepackage{titlesec}
\usepackage{fancyhdr}
\usepackage{enumitem}
\usepackage{pifont}
\usepackage{array}
\usepackage{colortbl}
\usepackage{hyperref}

% ------------------- CHARTE QIMMA -------------------
\definecolor{qteal}{HTML}{0d9488}
\definecolor{qtealdark}{HTML}{134e4a}
\definecolor{qamber}{HTML}{fbbf24}
\definecolor{qambertext}{HTML}{92400e}
\definecolor{ink}{HTML}{1f2937}
\definecolor{inkgray}{HTML}{6b7280}
\definecolor{tealpale}{HTML}{e6f5f3}
\colorlet{colDef}{qteal}
\definecolor{colProp}{HTML}{0f766e}
\definecolor{colEx}{HTML}{b45309}
\definecolor{colRem}{HTML}{9d174d}
\color{ink}

\graphicspath{{../../../../assets/}}
\newcommand{\logofile}{../../../../assets/qimma-logo.pdf}

% ------------------- METADONNEES -------------------
\newcommand{\matiere}{Mathématiques}
\newcommand{\niveau}{2\up{ème} Bac -- Sciences Mathématiques}
\newcommand{\chapitretitre}{Calcul intégral}

% ------------------- EN-TETE / PIED DE PAGE -------------------
\pagestyle{fancy}
\fancyhf{}
\renewcommand{\headrulewidth}{0pt}
\renewcommand{\footrulewidth}{0.5pt}
\fancyfoot[L]{\raisebox{-0.15cm}{\includegraphics[height=0.35cm]{\logofile}}~\small\sffamily\color{inkgray}Qimma}
\fancyfoot[C]{\small\sffamily\color{inkgray}\matiere\ -- \niveau}
\fancyfoot[R]{\small\sffamily\color{qtealdark}\thepage}

% ------------------- TITRES DE SECTION -------------------
\titleformat{\section}
  {\normalfont\sffamily\Large\bfseries\color{qtealdark}}
  {\tikz[baseline=-3.2pt]{\node[circle,fill=qteal,text=white,inner sep=0pt,minimum size=8mm]{\sffamily\bfseries\thesection};}}
  {10pt}{}
  [\vspace{-2mm}{\color{qamber}\titlerule[1.6pt]}\vspace{2mm}]
\titlespacing*{\section}{0pt}{16pt}{10pt}

\titleformat{\subsection}
  {\normalfont\sffamily\large\bfseries\color{qtealdark}}
  {\color{qteal}\ding{212}}{6pt}{}
\titlespacing*{\subsection}{0pt}{12pt}{6pt}

% ------------------- BOITES PEDAGOGIQUES -------------------
\newtcolorbox{fiche}[3][]{
  enhanced, breakable,
  colback=white, colframe=#2,
  boxrule=1pt, arc=2mm,
  top=7mm, left=4mm, right=4mm, bottom=3mm,
  drop shadow={black!25, shadow xshift=1mm, shadow yshift=-1mm},
  attach boxed title to top left={xshift=4mm, yshift=-3mm, yshifttext=-1mm},
  boxed title style={colback=#2, arc=1.5mm, boxrule=0pt, left=2mm, right=2mm},
  coltitle=white, fonttitle=\sffamily\bfseries,
  title={#3}, #1
}
\newenvironment{definition}[1][Définition]
  {\begin{fiche}{colDef}{\ding{110}~~#1}}{\end{fiche}}
\newenvironment{propriete}[1][Propriété]
  {\begin{fiche}{colProp}{\ding{72}~~#1}}{\end{fiche}}
\newenvironment{exemple}[1][Exemple]
  {\begin{fiche}{colEx}{\ding{217}~~#1}}{\end{fiche}}
\newenvironment{remarque}[1][Remarque]
  {\begin{fiche}{colRem}{\ding{43}~~#1}}{\end{fiche}}

% Démonstration (hors boîte, style discret)
\newenvironment{demo}
  {\par\smallskip\noindent{\sffamily\bfseries\color{qtealdark}Démonstration.}\ \itshape}
  {\ \normalfont\hfill$\blacksquare$\par\medskip}

\hypersetup{colorlinks=true, linkcolor=qtealdark, urlcolor=qtealdark}

% ============================================================
\begin{document}

% ------------------- BANNIERE DE TITRE -------------------
\begin{tcolorbox}[
  enhanced, boxrule=0pt, arc=4mm,
  interior style={left color=qtealdark, right color=qteal},
  top=8mm, bottom=8mm, left=10mm, right=32mm,
  drop shadow={black!35, shadow xshift=1.5mm, shadow yshift=-1.5mm},
  before skip=0pt, after skip=9mm,
  overlay={
    \node[anchor=north west] at ([xshift=6mm,yshift=-6mm]frame.north west)
      {\includegraphics[height=1.25cm]{\logofile}};
    \node[anchor=north east, fill=qamber, text=qambertext, rounded corners=2pt,
          inner sep=4pt, font=\sffamily\bfseries\small] at ([xshift=-6mm,yshift=-6mm]frame.north east) {COURS};
  }
]
\hspace{1.6cm}\centering
{\sffamily\bfseries\color{white}\fontsize{23}{27}\selectfont \chapitretitre}\\[3mm]
{\sffamily\color{white!90}\large \matiere\ \textbullet\ \niveau}
\end{tcolorbox}

% ------------------- OBJECTIFS -------------------
\begin{tcolorbox}[
  enhanced, colback=tealpale, colframe=qteal, boxrule=0.8pt, arc=2mm,
  left=5mm, right=5mm, top=3mm, bottom=3mm,
  title={\sffamily\bfseries\color{qtealdark}\ding{51}~~Objectifs du chapitre},
  coltitle=qtealdark, colbacktitle=tealpale, boxrule=0.8pt,
  attach title to upper={\par\vspace{1mm}}
]
\begin{itemize}[leftmargin=6mm, label=\ding{212}, itemsep=1mm]
  \item Calculer une intégrale à l'aide d'une primitive et démontrer les propriétés fondamentales (linéarité, Chasles, positivité, ordre).
  \item Interpréter géométriquement l'intégrale d'une fonction positive comme une aire, et l'étendre au cas d'une fonction de signe quelconque.
  \item Maîtriser complètement la technique d'intégration par parties, y compris les cas nécessitant deux applications successives.
  \item Calculer la valeur moyenne d'une fonction et l'exploiter dans un contexte physique (vitesse moyenne, etc.).
  \item Calculer des aires entre courbes et des volumes de solides de révolution, avec tous les cas de figure (fonction de signe variable, deux courbes).
  \item Reconnaître une somme de Riemann et l'utiliser pour calculer la limite d'une suite définie par une somme.
  \item Étudier une fonction définie par une intégrale $F(x)=\int_a^x f(t)\,dt$ sans nécessairement expliciter de primitive.
  \item Étudier une suite définie par une intégrale $u_n=\int_a^bf_n(t)\,dt$ par la méthode de l'encadrement.
  \item Résoudre les exercices type examen national combinant primitives, IPP et interprétation géométrique.
\end{itemize}
\end{tcolorbox}

\vspace{4mm}

% ------------------- SECTION 1 -------------------
\section{Intégrale d'une fonction continue}

\begin{definition}
Soit $f$ continue sur un intervalle $I$, $F$ une primitive de $f$ sur $I$, et $a,b\in I$. L'\textbf{intégrale de $a$ à $b$ de $f$} est le réel
\[
  \int_a^b f(x)\,dx = \Bigl[F(x)\Bigr]_a^b = F(b)-F(a).
\]
\end{definition}

\begin{propriete}[Indépendance vis-à-vis de la primitive choisie]
La valeur de $\displaystyle\int_a^b f(x)\,dx$ ne dépend pas de la primitive $F$ choisie.
\end{propriete}

\begin{demo}
Si $F$ et $G$ sont deux primitives de $f$ sur $I$, alors $G=F+C$ pour une constante $C$ (théorème du chapitre précédent). Donc $G(b)-G(a) = \bigl(F(b)+C\bigr)-\bigl(F(a)+C\bigr) = F(b)-F(a)$ : le résultat est bien indépendant du choix de $F$.
\end{demo}

\begin{remarque}[Variable muette]
$\displaystyle\int_a^b f(x)\,dx = \int_a^b f(t)\,dt$ : la lettre utilisée pour la variable d'intégration n'a pas d'importance, elle « disparaît » après calcul.
\end{remarque}

\begin{propriete}[Interprétation géométrique]
Si $f$ est \textbf{continue et positive} sur $[a\,;b]$, alors $\displaystyle\int_a^b f(x)\,dx$ est l'\textbf{aire}, en unités d'aire (u.a.), du domaine délimité par $\mathcal{C}_f$, l'axe des abscisses et les droites $x=a$, $x=b$. L'unité d'aire est l'aire du rectangle construit sur les vecteurs de base du repère.
\end{propriete}

\begin{exemple}[Premier calcul]
\[
  \int_1^2\left(3x^2-2x\right)dx = \Bigl[x^3-x^2\Bigr]_1^2 = (8-4)-(1-1) = 4.
\]
\end{exemple}

% ------------------- SECTION 2 -------------------
\section{Propriétés de l'intégrale}

\begin{propriete}[Propriétés fondamentales]
Pour $f,g$ continues sur $I$, $a,b,c\in I$, $\lambda\in\mathbb{R}$ :
\begin{itemize}[leftmargin=6mm, itemsep=0.5mm]
  \item $\displaystyle\int_a^a f(x)\,dx = 0$ ; \quad $\displaystyle\int_b^a f(x)\,dx = -\int_a^b f(x)\,dx$ ;
  \item \textbf{Linéarité} : $\displaystyle\int_a^b\bigl(\lambda f+g\bigr)(x)\,dx = \lambda\int_a^b f(x)\,dx+\int_a^b g(x)\,dx$ ;
  \item \textbf{Relation de Chasles} : $\displaystyle\int_a^b f(x)\,dx = \int_a^c f(x)\,dx+\int_c^b f(x)\,dx$ ;
  \item \textbf{Positivité} : si $a\leq b$ et $f\geq0$ sur $[a\,;b]$, alors $\displaystyle\int_a^b f(x)\,dx\geq0$ ;
  \item \textbf{Croissance/ordre} : si $a\leq b$ et $f\leq g$ sur $[a\,;b]$, alors $\displaystyle\int_a^b f(x)\,dx\leq\int_a^b g(x)\,dx$ ;
  \item \textbf{Inégalité de la moyenne} : si $m\leq f(x)\leq M$ sur $[a\,;b]$ ($a\leq b$), alors $m(b-a)\leq\displaystyle\int_a^b f(x)\,dx\leq M(b-a)$.
\end{itemize}
\end{propriete}

\begin{demo}
Toutes ces propriétés se démontrent directement à partir de la définition avec une primitive. Par exemple pour la \textbf{linéarité} : si $F$ est une primitive de $f$ et $G$ une primitive de $g$, alors $\lambda F+G$ est une primitive de $\lambda f+g$ (dérivée d'une somme et d'un produit par une constante), donc
\[
  \int_a^b(\lambda f+g) = \bigl[\lambda F+G\bigr]_a^b = \lambda\bigl(F(b)-F(a)\bigr)+\bigl(G(b)-G(a)\bigr) = \lambda\int_a^b f+\int_a^b g.
\]
Pour la \textbf{positivité} : $f\geq0$ sur $[a\,;b]$ signifie que $F$ est croissante sur $[a\,;b]$ (car $F'=f\geq0$), donc $F(b)\geq F(a)$, soit $\int_a^bf=F(b)-F(a)\geq0$.

\smallskip
La \textbf{croissance} se déduit de la positivité appliquée à $g-f\geq0$, et Chasles se vérifie en écrivant $F(b)-F(a) = \bigl(F(c)-F(a)\bigr)+\bigl(F(b)-F(c)\bigr)$.
\end{demo}

\begin{definition}[Valeur moyenne]
La \textbf{valeur moyenne} de $f$ continue sur $[a\,;b]$ ($a<b$) est
\[
  \mu = \frac{1}{b-a}\int_a^b f(x)\,dx.
\]
\end{definition}

\begin{exemple}[Encadrement sans calcul explicite]
Encadrons $I=\displaystyle\int_0^1\frac{1}{1+x^2}\,dx$ sans la calculer.

\smallskip
Pour $x\in[0\,;1]$ : $1\leq1+x^2\leq2$, donc $\frac12\leq\frac{1}{1+x^2}\leq1$. Par l'inégalité de la moyenne (avec $b-a=1$) : $\frac12\leq I\leq1$.
\end{exemple}

\begin{exemple}[Signe et fonction de signe variable]
Calculons $\displaystyle\int_{-1}^{2}x\,dx$ et interprétons.

\smallskip
$\displaystyle\int_{-1}^{2}x\,dx = \left[\frac{x^2}{2}\right]_{-1}^{2} = 2-\frac12 = \frac32$. Ce nombre est la \textbf{somme algébrique} des aires : positive sur $[0\,;2]$ (au-dessus de l'axe), négative sur $[-1\,;0]$ (en dessous), le résultat global n'est \textbf{pas} l'aire géométrique totale du domaine coloré (qui vaudrait $\int_{-1}^0(-x)dx+\int_0^2xdx = \frac12+2=\frac52$).
\end{exemple}

% ------------------- SECTION 3 -------------------
\section{Fonctions définies par une intégrale}

\begin{propriete}[Fonction intégrale -- primitive s'annulant en $a$]
Soit $f$ continue sur un intervalle $I$, et $a\in I$. La fonction
\[
  F : x \longmapsto \int_a^x f(t)\,dt
\]
est \textbf{dérivable} sur $I$, avec $F'=f$ et $F(a)=0$. Autrement dit, $F$ est \textbf{l'unique primitive} de $f$ sur $I$ qui s'annule en $a$.
\end{propriete}

\begin{demo}
$f$ est continue sur $I$, donc elle admet une primitive $G$ sur $I$ (théorème d'existence, chapitre précédent). Pour tout $x\in I$ :
\[
  F(x) = \int_a^x f(t)\,dt = G(x)-G(a)
\]
(par définition de l'intégrale à l'aide d'une primitive). Comme $G$ est dérivable de dérivée $f$ et que $G(a)$ est une constante, $F$ est dérivable et $F'(x)=G'(x)=f(x)$. De plus $F(a)=G(a)-G(a)=0$.
\end{demo}

\begin{remarque}[L'intérêt de ce théorème]
Il permet d'étudier une fonction $F$ définie par une intégrale (sens de variation, signe, limites) \textbf{sans jamais expliciter de primitive} : il suffit de connaître le signe de $f=F'$. C'est essentiel lorsque $f$ n'a pas de primitive exprimable avec les fonctions usuelles.
\end{remarque}

\begin{exemple}[Étudier $F$ sans expliciter de primitive]
Soit $F(x)=\displaystyle\int_0^xe^{-t^2}\,dt$ pour $x\in\mathbb{R}$. La fonction $t\mapsto e^{-t^2}$ n'admet \textbf{pas} de primitive exprimable à l'aide des fonctions usuelles (polynômes, $\ln$, $\exp$, fonctions trigonométriques) -- pourtant $F$ est parfaitement bien définie et étudiable grâce au théorème précédent.

\smallskip
$F$ est dérivable sur $\mathbb{R}$ (car $t\mapsto e^{-t^2}$ est continue sur $\mathbb{R}$), et $F'(x)=e^{-x^2}>0$ pour tout $x\in\mathbb{R}$ : $F$ est \textbf{strictement croissante} sur $\mathbb{R}$, avec $F(0)=0$ (car $\int_0^0=0$). Donc $F(x)<0$ pour $x<0$ et $F(x)>0$ pour $x>0$.
\end{exemple}

\begin{propriete}[Cas composé : $x\mapsto\int_a^{u(x)}f(t)\,dt$]
Soit $f$ continue sur $I$, $a\in I$, et $u$ dérivable sur un intervalle $J$ avec $u(J)\subset I$. Alors la fonction
\[
  H : x \longmapsto \int_a^{u(x)}f(t)\,dt
\]
est dérivable sur $J$, et
\[
  H'(x) = u'(x)\times f\bigl(u(x)\bigr).
\]
\end{propriete}

\begin{demo}
Posons $F(x)=\int_a^xf(t)\,dt$ (fonction du théorème précédent, $F'=f$). Alors $H=F\circ u$, composée de $u$ (dérivable) et $F$ (dérivable, de dérivée $f$). Par la formule de dérivation d'une composée :
\[
  H'(x) = u'(x)\times F'\bigl(u(x)\bigr) = u'(x)\times f\bigl(u(x)\bigr).
\]
\end{demo}

\begin{exemple}[Étudier une fonction composée définie par une intégrale]
Soit $H(x)=\displaystyle\int_0^{x^2}e^{-t^2}\,dt$ pour $x\in\mathbb{R}$. Étudions les variations de $H$.

\smallskip
Ici $u(x)=x^2$ (dérivable, $u'(x)=2x$) et $f(t)=e^{-t^2}$ (continue sur $\mathbb{R}$). Par la formule précédente :
\[
  H'(x) = u'(x)\times f\bigl(u(x)\bigr) = 2x\times e^{-x^4}.
\]
Comme $e^{-x^4}>0$ pour tout $x$, le signe de $H'$ est celui de $2x$ : $H$ est décroissante sur $]-\infty\,;0]$, croissante sur $[0\,;+\infty[$, avec $H(0)=\int_0^0=0$ (minimum).
\end{exemple}

% ------------------- SECTION 4 -------------------
\section{Intégration par parties}

\begin{propriete}[Formule d'intégration par parties]
Si $u$ et $v$ sont dérivables sur $[a\,;b]$, de dérivées $u'$ et $v'$ continues :
\[
  \int_a^b u(x)v'(x)\,dx = \Bigl[u(x)v(x)\Bigr]_a^b - \int_a^b u'(x)v(x)\,dx.
\]
\end{propriete}

\begin{demo}
$(uv)' = u'v+uv'$, donc $uv$ est une primitive de $u'v+uv'$. En intégrant cette égalité entre $a$ et $b$ :
\[
  \Bigl[u(x)v(x)\Bigr]_a^b = \int_a^b u'(x)v(x)\,dx + \int_a^b u(x)v'(x)\,dx,
\]
ce qui donne, en isolant le terme voulu, la formule annoncée.
\end{demo}

\begin{remarque}[Bien choisir $u$ et $v'$]
On choisit $u$ = la fonction qui se \textbf{simplifie} en dérivant (polynôme, $\ln x$) et $v'$ = celle qu'on sait intégrer facilement ($e^x$, $\cos x$, $\sin x$, polynôme). Cas typiques : $\int xe^xdx$, $\int x\cos xdx$ (poser $u=x$) ; $\int x\ln xdx$, $\int\ln xdx$ (poser $u=\ln x$, $v'=1$).
\end{remarque}

\begin{exemple}[IPP classique : $xe^x$]
Calculons $I=\displaystyle\int_0^1xe^x\,dx$. Posons $u(x)=x$ ($u'=1$), $v'(x)=e^x$ ($v=e^x$) :
\[
  I = \Bigl[xe^x\Bigr]_0^1-\int_0^1e^x\,dx = e-\Bigl[e^x\Bigr]_0^1 = e-(e-1) = 1.
\]
\end{exemple}

\begin{exemple}[IPP avec le logarithme]
Calculons $J=\displaystyle\int_1^e\ln x\,dx$. Posons $u(x)=\ln x$ ($u'=\frac1x$), $v'(x)=1$ ($v=x$) :
\[
  J = \Bigl[x\ln x\Bigr]_1^e-\int_1^e x\cdot\frac1x\,dx = e\ln e-1\ln1-\int_1^e1\,dx = e-0-(e-1) = 1.
\]
\end{exemple}

\begin{exemple}[IPP double : $x^2e^x$]
Calculons $K=\displaystyle\int_0^1x^2e^x\,dx$ (nécessite deux IPP successives).

\smallskip
\textbf{Première IPP} : $u=x^2$ ($u'=2x$), $v'=e^x$ ($v=e^x$) :
\[
  K = \Bigl[x^2e^x\Bigr]_0^1-\int_0^12xe^x\,dx = e - 2\int_0^1xe^x\,dx.
\]
\textbf{Deuxième IPP} (reprend l'exemple précédent) : $\displaystyle\int_0^1xe^x\,dx=1$.

\smallskip
\textbf{Conclusion} : $K = e-2\times1 = e-2$.
\end{exemple}

% ------------------- SECTION 5 -------------------
\section{Applications : aires et volumes}

\begin{propriete}[Calcul d'aires]
\begin{itemize}[leftmargin=6mm, itemsep=0.5mm]
  \item Si $f\geq0$ sur $[a\,;b]$ : aire $=\displaystyle\int_a^bf(x)\,dx$ (u.a.).
  \item Si $f\leq0$ sur $[a\,;b]$ : l'aire entre la courbe et l'axe des abscisses vaut $-\displaystyle\int_a^bf(x)\,dx$.
  \item Si $f$ change de signe : on \textbf{découpe} $[a\,;b]$ selon le signe de $f$ (relation de Chasles) et on somme les \textbf{valeurs absolues} des intégrales sur chaque sous-intervalle.
  \item \textbf{Aire entre deux courbes} : si $f\geq g$ sur $[a\,;b]$, l'aire du domaine compris entre les deux courbes vaut
  \[
    \mathcal{A} = \int_a^b\bigl[f(x)-g(x)\bigr]\,dx \quad\text{(u.a.)}.
  \]
\end{itemize}
\end{propriete}

\begin{propriete}[Volume d'un solide de révolution]
Le solide engendré par la rotation autour de l'axe $(Ox)$ du domaine sous la courbe de $f$ sur $[a\,;b]$ a pour volume :
\[
  V = \pi\int_a^b\bigl[f(x)\bigr]^2\,dx \quad\text{(u.v.)}.
\]
\end{propriete}

\begin{exemple}[Aire d'une fonction de signe variable]
Calculons l'aire du domaine compris entre $\mathcal{C}_f$, l'axe $(Ox)$, et les droites $x=0$, $x=2$, où $f(x)=x-1$.

\smallskip
$f$ s'annule en $x=1$ : $f\leq0$ sur $[0\,;1]$, $f\geq0$ sur $[1\,;2]$.
\[
  \mathcal{A} = -\int_0^1(x-1)\,dx + \int_1^2(x-1)\,dx
  = -\left[\frac{x^2}2-x\right]_0^1 + \left[\frac{x^2}2-x\right]_1^2
  = \frac12+\frac12 = 1 \text{ u.a.}
\]
\end{exemple}

\begin{exemple}[Aire entre deux courbes]
Calculons l'aire du domaine compris entre $f(x)=x$ et $g(x)=x^2$ pour $x\in[0\,;1]$.

\smallskip
Sur $[0\,;1]$ : $x\geq x^2$ (car $x-x^2=x(1-x)\geq0$). Donc
\[
  \mathcal{A} = \int_0^1(x-x^2)\,dx = \left[\frac{x^2}2-\frac{x^3}3\right]_0^1 = \frac12-\frac13 = \frac16 \text{ u.a.}
\]
\end{exemple}

\begin{exemple}[Volume : cône de révolution]
Retrouvons le volume du cône engendré par la rotation de $y=\frac rhx$ ($x\in[0\,;h]$) autour de $(Ox)$ :
\[
  V = \pi\int_0^h\frac{r^2}{h^2}x^2\,dx = \frac{\pi r^2}{h^2}\left[\frac{x^3}3\right]_0^h = \frac{\pi r^2}{h^2}\cdot\frac{h^3}3 = \frac13\pi r^2h.
\]
On retrouve bien la formule connue du volume du cône.
\end{exemple}

% ------------------- SECTION 6 -------------------
\section{Sommes de Riemann : limites de sommes liées à une intégrale}

\begin{propriete}[Sommes de Riemann -- admis]
Si $f$ est continue sur $[a\,;b]$, alors
\[
  \lim_{n\to+\infty}\frac{b-a}{n}\sum_{k=0}^{n-1}f\left(a+k\,\frac{b-a}{n}\right) = \int_a^bf(x)\,dx.
\]
Cas particulier fondamental ($a=0$, $b=1$) :
\[
  \lim_{n\to+\infty}\frac1n\sum_{k=0}^{n-1}f\left(\frac kn\right) = \int_0^1f(x)\,dx.
\]
De façon équivalente (même limite, indice décalé de $1$) :
\[
  \lim_{n\to+\infty}\frac1n\sum_{k=1}^{n}f\left(\frac kn\right) = \int_0^1f(x)\,dx.
\]
\end{propriete}

\begin{remarque}[Interprétation]
$\dfrac{b-a}{n}\displaystyle\sum_{k=0}^{n-1}f\left(a+k\frac{b-a}{n}\right)$ est l'aire d'une réunion de $n$ rectangles de largeur $\frac{b-a}{n}$, approchant l'aire sous la courbe de $f$ : quand $n\to+\infty$, cette approximation \textbf{converge vers l'intégrale}. C'est ce résultat qui justifie, en amont, l'interprétation de l'intégrale comme une aire.
\end{remarque}

\begin{exemple}[Calculer une limite de somme via une intégrale]
Calculons $\displaystyle\lim_{n\to+\infty}\frac1n\sum_{k=1}^{n}\sqrt{\frac kn}$.

\smallskip
On reconnaît une somme de Riemann associée à $f(x)=\sqrt x$ sur $[0\,;1]$ :
\[
  \frac1n\sum_{k=1}^n\sqrt{\frac kn} = \frac1n\sum_{k=1}^nf\left(\frac kn\right) \xrightarrow[n\to+\infty]{} \int_0^1\sqrt x\,dx = \left[\frac23x^{3/2}\right]_0^1 = \frac23.
\]
\end{exemple}

\begin{exemple}[Somme avec fonction exponentielle]
Calculons $\displaystyle\lim_{n\to+\infty}\frac1n\sum_{k=1}^{n}e^{k/n}$.

\smallskip
On reconnaît une somme de Riemann associée à $f(x)=e^x$ sur $[0\,;1]$ :
\[
  \frac1n\sum_{k=1}^ne^{k/n} = \frac1n\sum_{k=1}^nf\left(\frac kn\right) \xrightarrow[n\to+\infty]{} \int_0^1e^x\,dx = \Bigl[e^x\Bigr]_0^1 = e-1.
\]
\end{exemple}

\subsection{Suites définies par une intégrale}

\begin{remarque}[Méthode de l'encadrement]
Pour étudier une suite du type $u_n=\displaystyle\int_a^b f_n(t)\,dt$, où $f_n$ dépend de $n$ (paramètre \textbf{dans l'intégrande}, à ne pas confondre avec les fonctions $F$ et $H$ des paragraphes précédents où $n$ ou $x$ intervenaient dans les \textbf{bornes}), on encadre $f_n(t)$ pour $t\in[a\,;b]$, puis on \textbf{intègre l'encadrement} (croissance de l'intégrale) pour encadrer $u_n$, et on conclut par le théorème des gendarmes.
\end{remarque}

\begin{exemple}[Suite définie par une intégrale]
Soit $u_n=\displaystyle\int_0^1\frac{x^n}{1+x}\,dx$ pour $n\in\mathbb{N}$. Montrons que $u_n\to0$.

\smallskip
Pour $x\in[0\,;1]$ : $1\leq1+x\leq2$, donc, en passant à l'inverse puis en multipliant par $x^n\geq0$ :
\[
  \frac{x^n}2 \leq \frac{x^n}{1+x} \leq x^n.
\]
En intégrant sur $[0\,;1]$ (croissance de l'intégrale) :
\[
  \int_0^1\frac{x^n}2\,dx \leq u_n \leq \int_0^1x^n\,dx,
  \qquad\text{soit}\qquad
  \frac{1}{2(n+1)} \leq u_n \leq \frac{1}{n+1}.
\]
Comme $\dfrac{1}{2(n+1)}\to0$ et $\dfrac{1}{n+1}\to0$ : par le théorème des gendarmes, $u_n\to0$.
\end{exemple}

% ------------------- SECTION 7 -------------------
\section{Exercices corrigés -- niveau examen national SM}

\begin{exemple}[Exercice 1 -- intégrale avec fonction composée $\ln$]
Calculer $\displaystyle\int_1^e\frac{(\ln x)^2}{x}\,dx$.

\smallskip
\textbf{Corrigé.} On reconnaît $u'u^2$ avec $u=\ln x$ (car $u'=\frac1x$) : une primitive est $\frac{(\ln x)^3}{3}$. Donc
\[
  \int_1^e\frac{(\ln x)^2}{x}\,dx = \left[\frac{(\ln x)^3}{3}\right]_1^e = \frac{1}{3}-0 = \frac13.
\]
\end{exemple}

\begin{exemple}[Exercice 2 -- aire limitée par une exponentielle]
Calculer l'aire du domaine délimité par $\mathcal{C}_f$ ($f(x)=e^{-x}$), l'axe $(Ox)$, l'axe $(Oy)$ et la droite $x=\ln2$.

\smallskip
\textbf{Corrigé.} $f\geq0$ sur $[0\,;\ln2]$, donc
\[
  \mathcal{A} = \int_0^{\ln2}e^{-x}\,dx = \Bigl[-e^{-x}\Bigr]_0^{\ln2} = -e^{-\ln2}-(-1) = -\frac12+1 = \frac12 \text{ u.a.}
\]
\end{exemple}

\begin{exemple}[Exercice 3 -- IPP avec borne variable et fonction $F$]
Soit $F(x)=\displaystyle\int_0^x t\,e^{-t}\,dt$ pour $x\geq0$. Calculer $F(x)$ explicitement, puis étudier $\lim\limits_{x\to+\infty}F(x)$.

\smallskip
\textbf{Corrigé.} IPP avec $u=t$ ($u'=1$), $v'=e^{-t}$ ($v=-e^{-t}$) :
\[
  \int_0^xte^{-t}\,dt = \Bigl[-te^{-t}\Bigr]_0^x+\int_0^xe^{-t}\,dt = -xe^{-x}+\Bigl[-e^{-t}\Bigr]_0^x = -xe^{-x}-e^{-x}+1.
\]
Donc $F(x)=1-(x+1)e^{-x}$. Quand $x\to+\infty$ : $(x+1)e^{-x}\to0$ (croissance comparée), donc $F(x)\to1$.
\end{exemple}

\begin{exemple}[Exercice 4 -- volume avec fonction $\ln$]
Calculer le volume engendré par la rotation autour de $(Ox)$ du domaine sous la courbe de $f(x)=\sqrt{\ln x}$ pour $x\in[1\,;e]$.

\smallskip
\textbf{Corrigé.}
\[
  V = \pi\int_1^e\bigl(\sqrt{\ln x}\bigr)^2\,dx = \pi\int_1^e\ln x\,dx.
\]
Or $\displaystyle\int_1^e\ln x\,dx = 1$ (calculé à l'exemple de la section 3). Donc $V=\pi$ (u.v.).
\end{exemple}

\begin{exemple}[Exercice 5 -- comparaison d'intégrales via une inégalité]
Sachant que $\ln x \leq x-1$ pour tout $x>0$, en déduire un encadrement de $\displaystyle\int_1^2\ln x\,dx$ sans calculer explicitement de primitive du membre de gauche autrement.

\smallskip
\textbf{Corrigé.} Pour $x\in[1\,;2]$ : $\ln x\leq x-1$. Par croissance de l'intégrale :
\[
  \int_1^2\ln x\,dx \leq \int_1^2(x-1)\,dx = \left[\frac{x^2}2-x\right]_1^2 = (2-2)-\left(\frac12-1\right) = \frac12.
\]
\emph{Vérification par calcul exact} (IPP comme précédemment) : $\int_1^2\ln x\,dx = \bigl[x\ln x-x\bigr]_1^2 = (2\ln2-2)-(0-1) = 2\ln2-1\approx0{,}386 \leq0{,}5$. $\checkmark$
\end{exemple}

% ------------------- SECTION 8 -------------------
\section{Méthodes et pièges : la boîte à outils de l'examen}

\subsection{Ce qu'on te demande, ce que tu fais}

\begin{center}
\renewcommand{\arraystretch}{1.45}
\sffamily\small
\begin{tabular}{>{\raggedright\arraybackslash}p{7.2cm} >{\raggedright\arraybackslash}p{8cm}}
\rowcolor{qtealdark} \color{white}\bfseries Ce qu'on te demande & \color{white}\bfseries Ton réflexe \\
\rowcolor{tealpale} Calculer une intégrale directe & Trouver une primitive $F$, calculer $F(b)-F(a)$. \\
Produit de deux fonctions de nature différente & Intégration par parties : $u$ se simplifie, $v'$ s'intègre facilement. \\
\rowcolor{tealpale} Puissance de $x$ multipliée par $e^x$ ou $\cos x$ & IPP répétée (parfois deux fois), en réduisant le degré à chaque étape. \\
Encadrer une intégrale sans la calculer & Encadrer $f$ sur $[a\,;b]$ puis intégrer les inégalités (croissance de l'intégrale). \\
\rowcolor{tealpale} Aire sous une courbe qui change de signe & Découper selon le signe de $f$, sommer les \textbf{valeurs absolues}. \\
Aire entre deux courbes & Étudier le signe de $f-g$, intégrer $|f-g|$ (ou $f-g$ si $f\geq g$ constaté). \\
\rowcolor{tealpale} Volume de révolution & $V=\pi\int_a^b f(x)^2\,dx$, ne pas oublier le carré. \\
Étudier $F(x)=\int_a^x g(t)\,dt$ & $F'=g$ (dérivable, primitive de $g$ nulle en $a$) ; variations données par le signe de $g$, même sans primitive explicite. \\
\rowcolor{tealpale} Étudier $H(x)=\int_a^{u(x)} g(t)\,dt$ & $H'(x)=u'(x)\times g(u(x))$ (composée avec la fonction intégrale). \\
\rowcolor{tealpale} Limite d'une somme $\frac1n\sum f\left(\frac kn\right)$ & Reconnaître une somme de Riemann, identifier $f$, conclure $\to\int_0^1f(x)\,dx$. \\
Limite d'une suite $u_n=\int_a^bf_n(t)\,dt$ & Encadrer $f_n(t)$ sur $[a\,;b]$, intégrer l'encadrement, conclure par les gendarmes. \\
\end{tabular}
\end{center}

\subsection{Les pièges qui coûtent des points}

\begin{remarque}[Erreurs relevées chaque année par les correcteurs]
\begin{itemize}[leftmargin=6mm, itemsep=0.5mm]
  \item Confondre \textbf{intégrale} (nombre algébrique, peut être négatif) et \textbf{aire} (toujours positive) quand $f$ change de signe.
  \item Oublier le \textbf{carré} dans la formule du volume de révolution $V=\pi\int f^2$.
  \item Se tromper dans le choix de $u$ et $v'$ pour l'IPP (résultat : intégrale plus compliquée qu'au départ -- signe qu'il faut inverser le choix).
  \item Oublier un signe « $-$ » en primitivant $\cos$, $\sin$ ou lors du calcul de $-e^{-x}$.
  \item Utiliser l'inégalité de la moyenne avec $a>b$ sans adapter les inégalités (elle suppose $a\leq b$).
  \item Pour l'aire entre deux courbes : oublier de vérifier le signe de $f-g$ sur \textbf{tout} l'intervalle (il peut changer, nécessitant un découpage).
  \item Pour $F(x)=\int_a^xg(t)dt$ : chercher à tout prix une primitive explicite de $g$ alors que le signe de $g=F'$ suffit pour répondre (variations, signe de $F$).
\end{itemize}
\end{remarque}

% ------------------- RESUME -------------------
\vspace{4mm}
\begin{tcolorbox}[
  enhanced, boxrule=0pt, arc=2mm,
  interior style={left color=qtealdark, right color=qteal},
  left=6mm, right=6mm, top=4mm, bottom=4mm,
  fontupper=\color{white}\sffamily,
  title={\sffamily\bfseries\color{white}\ding{72}~~À retenir},
  colbacktitle=qtealdark, coltitle=white, boxrule=0pt
]
\begin{itemize}[leftmargin=6mm, label={\color{qamber}\ding{212}}, itemsep=1mm]
  \item $\int_a^bf(x)\,dx=F(b)-F(a)$, indépendant de la primitive choisie ; interprétation aire si $f\geq0$.
  \item Linéarité, Chasles, positivité, croissance, inégalité de la moyenne : les cinq outils démontrés à partir d'une primitive.
  \item IPP : $\int_a^buv'=[uv]_a^b-\int_a^bu'v$ ; parfois répétée pour les puissances de $x$.
  \item Aire signe variable : découper + valeurs absolues ; aire entre courbes : $\int(f-g)$ si $f\geq g$ ; volume : $V=\pi\int f^2$.
  \item Valeur moyenne : $\mu=\frac{1}{b-a}\int_a^bf$.
  \item Sommes de Riemann : $\frac1n\sum_{k=1}^nf\left(\frac kn\right) \to \int_0^1f(x)\,dx$.
  \item $F(x)=\int_a^xf(t)\,dt$ est dérivable, $F'=f$, $F(a)=0$ : on l'étudie via le signe de $f$, même sans primitive explicite.
  \item Cas composé : $H(x)=\int_a^{u(x)}f(t)\,dt$ est dérivable, $H'(x)=u'(x)\,f(u(x))$.
  \item Suite $u_n=\int_a^bf_n(t)\,dt$ : encadrer $f_n$ sur $[a\,;b]$, intégrer, conclure par les gendarmes.
\end{itemize}
\end{tcolorbox}

\end{document}
